\documentclass[11pt]{article}

\usepackage[margin=1.05in]{geometry}
\usepackage{amsmath,amssymb,amsthm,mathtools}
\usepackage{booktabs}
\usepackage{microtype}
\usepackage{xcolor}
\usepackage[colorlinks=true,linkcolor=blue!45!black,citecolor=blue!45!black,urlcolor=blue!45!black]{hyperref}

\newtheorem{theorem}{Theorem}
\theoremstyle{remark}
\newtheorem{remark}{Remark}

\newcommand{\sight}{\mathsf{sight}}
\newcommand{\los}{\mathsf{los}}
\newcommand{\HtoC}{\mathsf{H2C}}
\newcommand{\sha}{\mathrm{SHA\text{-}256}}
\newcommand{\seed}{\mathrm{seed}}
\newcommand{\pseed}{\seed^{\mathrm{path}}}
\newcommand{\petseed}{\seed^{\mathrm{pet}}}
\newcommand{\master}{\mathrm{master}}
\newcommand{\concat}{\,\|\,}

\title{Fog of War without Zero-Knowledge Proofs:\\
A Blinded Mutual-Sighting Test for Trustless Game Channels}

\author{Andrew Colosimo \quad Roy Crombleholme \quad Andrew Gore\\[3pt]
Konstantin Gorskov \quad Daniel Kraft\\[6pt]
\normalsize Xaya Developers}

\date{July 31, 2026}

\begin{document}
\maketitle

\begin{abstract}
Replicated deterministic state machines, blockchains and the state channels
built on them, make hidden information awkward: every replica must be able to
verify every transition, yet fog of war requires that positions stay secret.
Commit and reveal hides positions but destroys \emph{encounters}: two players
walk past each other in the dark because neither may learn that the other is
near. We present a \emph{blinded mutual-sighting test}. After both players have
hash-committed their moves for a round, and before either move is disclosed, the
two clients run a Diffie--Hellman-blinded private membership test, the
Meadows~1986 and Huberman--Franklin--Hogg~1999 lineage cut down from set
intersection to a single membership bit evaluated in both directions. Both
parties learn exactly one shared bit, ``we are in sight'', and nothing else.
Three design choices adapt the classical protocol to a consensus setting.
First, every blinding scalar and pad element is derived deterministically from
per-epoch secrets the player is already bound to reveal, which makes the honest
transcript recomputable by an after-the-fact audit --- and any reveal that
publishes a secret closes its epoch on the spot: the referee restarts the
player's commitment chains and the client re-seeds behind them, so a disclosed
secret never governs a live round. Second, no conviction ever trusts the
exchange itself: cheating is adjudicated by recomputing the public sight
predicate over positions fixed by commitments made before the round's
information existed. Third, the blinded query element each round publishes is
bound, retroactively, to the position committed beside it --- one sampled curve
check at the audit, plus a pay-on-accusation dispute over sealed epochs ---
pricing, at audit strength, the chosen-element oracle that unbound
private-membership exchanges leave open. Together these replace an in-round
zero-knowledge proof, which does not fit the execution budget of our referee,
with after-the-fact attribution. On the honest path the referee performs no
curve arithmetic --- its in-round work is refusing a published set that does
not rebuild the committed root, which is hashing; audits and disputes each buy
a fixed, measured handful of curve operations --- and the chain is touched only
to open, dispute and settle.
\end{abstract}

\section{Introduction}

Multiplayer games on public replicated state machines inherit a contradiction.
The chain, or the channel, exists so that every participant can verify every
transition; fog of war exists so that participants \emph{cannot} see part of the
state. The standard resolution, commit now and reveal later, solves only half of
the problem. A hash commitment hides a player's position perfectly, and perfect
hiding is precisely what makes the game unplayable: if neither player may learn
anything about the other's position before the endgame, then two players a few
tiles apart walk past each other in the dark, every round, for the whole match.
In our own testbed, removing all position disclosure produced long stretches of
mutual line of sight in which neither player was ever told anything, which is an
invisible non-game.

The requirement is therefore a conjunction:

\begin{enumerate}
\item \textbf{Privacy.} Nobody learns anything about your position unless they
are close to you, or unless you spend it deliberately, for instance on a
contested objective.
\item \textbf{Encounters.} Two players who \emph{are} close must both learn so,
in the same round, every round, with no exceptions.
\item \textbf{Adjudication.} Every rule above must remain enforceable by a
referee that never sees secret state, on a metered execution budget.
\end{enumerate}

The natural cryptographic answer to requirement~3 is a zero-knowledge proof
attached to every move. That answer is unavailable to us in-round: our referee
is a zero-import WebAssembly blob with a hard fuel budget per call, and a single
Groth16 verification costs more than the whole budget. This paper describes what
we built instead.

\paragraph{Contributions.}
\begin{enumerate}
\item A \emph{blinded mutual-sighting test} (\S\ref{sec:protocol}): a two-party
private membership test, run in both directions between move commitment and move
disclosure, that yields the one shared bit $\sight(a,b)$ and nothing else. The
construction is the classical Diffie--Hellman private-matching protocol of
Meadows and of Huberman, Franklin and Hogg~\cite{meadows86,hfh99}, specialised
from set$\,\cap\,$set to membership of a single element.
\item A \emph{determinism-for-auditability} discipline (\S\ref{sec:audit}): every
byte a party should send is a pure function of secrets it is already committed to
reveal, scoped to an \emph{epoch} that any reveal terminates
(\S\ref{sec:epochs}), so a judge can recompute any disputed transcript after the
fact and no reveal ever un-blinds a live round. This
converts ``prove your input was honest'', which is a proof system's job, into
``lying is non-repudiable evidence'', at negligible consensus cost.
\item \emph{Enforcement by geometry, not cryptography} (\S\ref{sec:audit}): no
conviction ever depends on the exchange. Suppressing a sighting, and disputing a
phantom one, are both adjudicated by recomputing the public sight predicate over
positions bound by hash commitments made before the round's information existed.
\item A \emph{binding for the blinded query element} (\S\ref{sec:qbind}): the one
input the exchange cannot validate in-round --- a responder able to tell what it
is being asked would defeat the instrument --- is made retroactively checkable
against the committed position instead, by a sampled check at the audit and a
dispute over sealed epochs in which every well-formed filing convicts one of
its two parties (a malformed filing is refused without a verdict). We measured
the oracle this prices and found it decisive rather than marginal
(\S\ref{sec:oracle}).
\end{enumerate}

\section{Setting and threat model}
\label{sec:setting}

\paragraph{Game channels.} The system runs on \emph{game channels}
\cite{kraft16} as deployed on the Xaya platform: two players play a turn-based
game off-chain, exchanging signed moves, with the blockchain as the fallback
arbiter. A single host process referees many different games by loading each
game's rules as a hash-pinned, zero-import WebAssembly blob; the blob is
fuel-metered per call and its state is size-capped. Both clients and every
potential arbiter run the identical blob, so ``the judge'' below means \emph{any}
party replaying the rules; honesty of the judge is by replication, not by
assumption.

\paragraph{Round structure.} Play proceeds in rounds on a fixed public map, in
our deployment a tile dungeon generated procedurally from a shared seed. Each
round:
\begin{enumerate}
\item \textbf{Commit.} Both players publish hash commitments to their move,
position and action, salted from a per-epoch path secret (\S\ref{sec:derive}).
The same move
carries a Merkle root over the flight the player is about to publish, with its
own blinded element beside it in the clear --- the ordering \S\ref{sec:audit}
depends on: every root is public before any flight exists on the wire. The
referee also appends that blinded element to an append-only Merkle frontier
carried in the player's state, the retention behind \S\ref{sec:qbind}.
\item \textbf{Sighting test.} The exchange of \S\ref{sec:protocol}, run as a
pass of its own: each player publishes its blinded set and its responses as
signed channel moves, and the referee refuses a published set that does not
rebuild the root committed in step 1. Responses are computed from the peer's
committed element read out of the published state, so no message is delivered
outside the channel.
\item \textbf{Disclose.} If the test said \emph{no}, a player publishes nothing
about its position. If it said \emph{yes}, both players are obliged to publish
their exact positions for the round, with an arrival proof, and on entering
contact the path connecting the position to the player's last judged one.
\end{enumerate}
Not every ending runs an audit. A match decided during play finishes on the
spot: by a kill (combat is legal only in contact, entering which
requires a verified full-path proof, so the fight's participants already stand
audited), by a forfeit, or by a conviction that leaves one player standing. The
round cap, and any player
demanding it, close instead with an \textbf{endgame audit}: each player reveals its
current secrets and its path back to its last \emph{epoch boundary}
(\S\ref{sec:epochs}), which the blob verifies against those rounds' commitments
and replays against the map. Rounds before a boundary were verified identically
at the moment their epoch closed, and the closure is pinned into the state by
preimage: no round of movement is verified twice, and none that a verdict leans
on escapes verification. Mid-match reveals --- first contact, dispute
evidence --- are epoch boundaries rather than exceptional events. The audit
predates the sighting test; \S\ref{sec:audit} extends it.

\paragraph{Adversary.} A malicious client controls all of its own messages: it
may send arbitrary bytes in the sighting exchange, stall, or disclose falsely. It
cannot break hash commitments or signatures, and it cannot dodge an audit once
one is owed except by forfeiting, since timeouts convert stalling into forfeit. Both
players' clients may be malicious simultaneously; the game between two colluding
cheaters is not our concern, but neither cheater may gain against an honest
opponent. The design rule throughout: \emph{consensus never trusts a claim}, and
everything that convicts is recomputed by the blob from committed data.

\section{The sight predicate}
\label{sec:predicate}

Positions are walkable tiles of the public map. For tiles $a,b$ let $d(a,b)$ be
the shortest walkable path distance and let $\los(a,b)$ be the game's
line-of-sight predicate, a symmetric visibility test on the tile grid. For a
fixed sight range $\rho$ the predicate jointly computed each round is
\begin{equation}
\sight(a,b) \;:=\; \bigl[\, d(a,b) \le \rho \,\bigr] \;\wedge\; \los(a,b).
\label{eq:sight}
\end{equation}
Here $d$ is symmetric by construction, and $\los$ was verified symmetric
exhaustively over all in-range walkable pairs of the test maps. Symmetry makes
``the'' shared bit well defined: writing
\begin{equation}
V(x) \;:=\; \{\, t \text{ walkable} : \sight(x,t) \,\},
\qquad
a \in V(b) \iff b \in V(a),
\label{eq:vset}
\end{equation}
with $x \in V(x)$ always, since distance 0 satisfies \eqref{eq:sight}. The
protocol pads every set to a fixed size $N$ chosen above the largest $|V|$ the
map admits (\S\ref{sec:derive}), so message size and work are constants
independent of the sender's surroundings.

The predicate must be symmetric for the construction to apply at all: asymmetric
sight ranges would make the two parties' bits legitimately disagree and the
shared-bit framing collapse. Consequently the game defines exactly one sight
range in consensus, and light sources are cosmetic.

\section{Cryptographic requirements}
\label{sec:prelim}

The protocol is not tied to a particular group. It needs three properties, and
only the third is easy to get wrong.

\begin{enumerate}
\item \textbf{Commuting blinding.} A group operation and scalars for which
$\alpha(\beta X) = \beta(\alpha X)$. Any prime-order group serves, and so does
RSA-style multiplicative blinding; nothing below depends on the choice.
\item \textbf{Hardness.} Distinguishing $\{\alpha P_t : t \in T_0\}$ from
$\{\alpha P_t : t \in T_1\}$ for candidate sets $T_0, T_1$, given the public
$P_t$, must be hard. In a prime-order elliptic-curve group this is a standard
decisional Diffie--Hellman problem --- for a passive observer of the
transcript. A \emph{responding} participant additionally obtains one blinded
evaluation $\alpha Q$ per round for a $Q$ it supplies, a one-query oracle for
$x \mapsto \alpha\,\HtoC(x)$, so claims against a receiver rest on assumptions
of the one-more or gap Diffie--Hellman kind~\cite{jl09,jl10};
\S\ref{sec:security} scopes which claim rests on which.
\item \textbf{Elements with unknown discrete logarithms.} The map
$t \mapsto P_t$ from a game position to a group element must be a hash-to-group
map whose outputs have discrete logarithms unknown to everyone, including the
party doing the hashing. This is load-bearing rather than stylistic, for the
reason given in \S\ref{sec:security}.
\end{enumerate}

\paragraph{Our instantiation.} We use Curve25519~\cite{bernstein06} with x-only
Montgomery arithmetic, the RFC~7748~\cite{rfc7748} ladder, and an x-only
Elligator~2~\cite{elligator} map with the cofactor cleared at construction so
that every point lies in the prime-order subgroup. The map is a single-map
encoding rather than the RFC~9380~\cite{rfc9380} suite: the protocol needs
determinism, unknown discrete logarithms and prime-order outputs; it does not
need the output to be indistinguishable from uniform, because every transmitted
element is additionally randomised by a per-round blinding scalar.
Interoperability is a non-goal, since every party runs the same pinned blob. The
choice matters only for cost, and the reference implementation~\cite{fowrepo}
carries the details, the known-answer vectors, and the argument for each
deviation.

\subsection{Derivations: nothing is random}
\label{sec:derive}

Every quantity in the exchange is derived rather than sampled. This is the part
that makes the audit of \S\ref{sec:audit} possible, so it is a protocol
requirement and not an implementation detail:
\begin{align}
\text{position elements (public):}\quad
& P_t = \HtoC\bigl(\mathtt{"T"} \concat \mathrm{enc}(t)\bigr),
\label{eq:tilepoint}\\[2pt]
\text{blinding scalar of player $i$, round $r$:}\quad
& \alpha_r^{(i)} = \sha\bigl(\mathtt{"P"} \concat \petseed_i \concat \mathrm{enc}(r)\bigr),
\label{eq:alpha}\\[2pt]
\text{pad dummies of player $i$, round $r$:}\quad
& D_{r,j}^{(i)} = \HtoC\bigl(\mathtt{"D"} \concat \petseed_i \concat \mathrm{enc}(r) \concat \mathrm{enc}(j)\bigr),
\label{eq:dummy}
\end{align}
with distinct domain-separation tags, the pad index running
$j = |V|,\dots,N-1$, and with $\alpha$ reduced into the scalar
range and the measure-zero zero case patched deterministically. Here
$\petseed_i$ is one half of a split pair of secrets, defined next.

The dummies of \eqref{eq:dummy} are deliberately \emph{round-dependent}, and we
know the property is load-bearing because an earlier revision lacked it. A
dispute filing publishes that round's $\alpha_r$ by design (\S\ref{sec:audit}),
and $\alpha_r$ is invertible modulo the group order, so an observer can divide
the filing round's published set by it and recover the raw pad elements. For the
filing round itself this is harmless --- a filing publishes the filer's exact
position anyway, which dominates $|V|$ --- but pad elements fixed per seed would
remain in service in every later round, recognisable there through the
chosen-element channel of \S\ref{sec:oracle}: each recovered element, replayed
as a query, yields a threshold bit about the current $|V|$ (the pad fills the
slots above $|V|$, so ``still in use'' bounds it by that element's index) ---
and with a fixed pad it would keep yielding that bit in every later round, a
standing oracle handle --- $|V|$ being the one quantity the fixed-size pad
exists to hide.
Folding the round into the
derivation confines any recovery to the round already disclosed. The cost is
that only the position elements of \eqref{eq:tilepoint} can be precomputed; the
pad is re-derived each round, the same $N$ hash-to-group maps the budget already
prices.

\paragraph{The two-secret split.} The secrets are two, and the distinction is
load-bearing. Player $i$ holds one master secret and derives
\begin{equation}
\pseed_i = \sha\bigl(\mathtt{"A"} \concat \master_i\bigr),
\qquad
\petseed_i = \sha\bigl(\mathtt{"B"} \concat \master_i\bigr),
\label{eq:split}
\end{equation}
the first salting every move commitment, the second --- the \emph{PET seed},
after the private equality test at the exchange's core --- feeding
\eqref{eq:alpha} and \eqref{eq:dummy} and nothing else. A full-path reveal
publishes both: the path seed, which is what the movement audit needs, and the
PET seed of the epoch that same transition closes --- the mint of
\S\ref{sec:epochs} is the transition's final step, so a disclosed PET seed
never governs a round still to be played. (The only reveal that carries no PET
seed is the empty one from a player whose restart was minted that very round.)
The one scalar a dispute's adjudication needs rides the filing itself,
unproven, and is pinned there by the adjudication's own equalities
(\S\ref{sec:audit}). Both sub-seeds are hash images, so publishing them reveals
nothing about the master and hence nothing about any later epoch. The reuse
claim above survives in scoped form: the honest transcript of an epoch is a
pure function of that epoch's two seeds, the committed positions and the public
map, and the player is bound to expose exactly those, in the reveal that ends
the epoch.

An earlier revision derived salts and blinding scalars from one seed, and the
coupling was a live defect rather than a stylistic blemish: the protocol's own
disclosure rule forces a full-path reveal at first contact, mid-match, after
which the one published seed made every past and future $\alpha_r$ computable
--- the player's published $Q_r$ then decodes to its exact tile by testing the
few thousand candidates, every remaining round. The split, with the epoch
machinery below, is what makes the privacy claims of \S\ref{sec:security} true
as stated rather than true only until first contact.

\subsection{Epochs: a published secret is a spent secret}
\label{sec:epochs}

Full-path reveals are not confined to the endgame. The disclosure rules force
one at first contact --- a player entering contact without a standing
proof must ship the connecting path, verified in that same call --- and a
sight dispute's filing always carries one as its evidence, even from a filer
already proven, because the adjudication anchors on the filer's judge-verified
tile for the disputed round. (The sealed-epoch dispute of \S\ref{sec:qbind}
carries none --- its evidence is entirely public --- and is not an epoch
boundary.) Each such reveal publishes the revealer's path seed, and a published
path seed opens every commitment it will ever salt, since a round's committed
tile and action span under $2^{19}$ candidates. A protocol that continued
play under a published seed would play the rest of the match face-up.

The referee therefore \emph{mints a restart} as the final step of any transition
that accepts a full-path reveal from a living player in an unfinished match. The
mint folds each of the player's commitment-chain heads into a fresh one with the
old head inside the preimage, so the verified prefix stays pinned and a restart
can never shed history; it records the just-verified position as the new epoch's
anchor, so a restart is not a licensed teleport; and the client re-seeds behind
it,
\begin{equation}
\master_i' \;=\; \sha\bigl(\mathtt{"E"} \concat \master_i \concat \mathrm{enc}(R)\bigr),
\label{eq:ratchet}
\end{equation}
re-deriving both sub-seeds from the new master by \eqref{eq:split}. The master
never appears on any wire, so the ratchet moves only forward: nothing published
in one epoch says anything about the secrets of the next.

The stretch between boundaries is an \emph{epoch}, and every audit is
epoch-relative: a reveal covers exactly the rounds since the revealer's
boundary and anchors to the boundary's judged position. Three consequences
matter below. A mid-match reveal costs exactly the epoch it closes --- which
that reveal published anyway --- and nothing after it. No round is ever
verified twice, so the total audit work over a match is linear in its length
rather than quadratic; a single audit call still pays for the epoch it opens,
which in the worst case is the whole match. And a dispute among
more than two players is no longer terminal for the table: the audit pass it
forces takes reveals only from the accused --- the filer's proof was parked at
its filing --- convicts on the evidence, and returns the survivors to play with
the spent epochs re-seeded. (The demand-an-immediate-end challenge move remains
terminal by design; a challenge the match could continue past would be a
purchasable snapshot of every player's position.)

\section{The protocol}
\label{sec:protocol}

Fix round $r$; both players have committed position and action for $r$. Write
$\alpha_A = \alpha_r^{(A)}$, and let $a, b$ be the committed positions. Before
any disclosure:

\paragraph{Flight 1 (both directions).}
\begin{equation}
A \to B:\quad
S_A = \mathrm{sort}\Bigl(
\bigl\{\, \alpha_A P_t : t \in V(a) \,\bigr\} \cup
\bigl\{\, \alpha_A D_{r,j}^{(A)} : j = |V(a)|,\dots,N-1 \,\bigr\}
\Bigr),
\qquad
Q_A = \alpha_A P_a,
\label{eq:flight1}
\end{equation}
and symmetrically $B \to A$ sends $(S_B, Q_B)$. Here $S$ is $N$ elements plus the
single element $Q$, and $\mathrm{sort}$ orders by encoding, a canonical order
independent of enumeration that reveals nothing about which entries are dummies.
Exactly $N{+}1$ hash-to-group maps and $N{+}1$ scalar multiplications are
performed whatever $|V|$ is, so neither length nor timing leaks the sender's
surroundings.

\paragraph{Flight 2 (both directions).}
\begin{equation}
A \to B:\quad R_A = \alpha_A Q_B \;(= \alpha_A \alpha_B P_b),
\qquad\qquad
B \to A:\quad R_B = \alpha_B Q_A \;(= \alpha_B \alpha_A P_a).
\label{eq:flight2}
\end{equation}

\paragraph{Local finish.}
\begin{equation}
\mathrm{bit}_A = \bigl[\, R_B \in \{\alpha_A s : s \in S_B\} \,\bigr],
\qquad
\mathrm{bit}_B = \bigl[\, R_A \in \{\alpha_B s : s \in S_A\} \,\bigr],
\label{eq:finish}
\end{equation}
each costing $N$ scalar multiplications and $N$ comparisons.

\begin{theorem}[Correctness]
\label{thm:correct}
Except with probability negligible in the hash length,
$\mathrm{bit}_A = \mathrm{bit}_B = \sight(a,b)$.
\end{theorem}

\begin{proof}[Proof sketch]
On the prime-order group, $\alpha_A(\alpha_B P_t) = \alpha_B(\alpha_A P_t)$, so
$R_B = \alpha_A\alpha_B P_a$ equals $\alpha_A(\alpha_B P_t)$ if and only if
$P_a = P_t$, that is if and only if $a \in V(b)$, unless two distinct
hash-to-group inputs collide, which is bounded by the birthday term over the
domain-separated inputs. Thus $\mathrm{bit}_A = [a \in V(b)]$ and symmetrically
$\mathrm{bit}_B = [b \in V(a)]$, and these coincide with $\sight(a,b)$ by
\eqref{eq:vset}. With x-only encodings $P$ and $-P$ are identified, but
hash-to-group outputs are themselves x-coordinates, so that coincidence is the
same collision event and is already counted.
\end{proof}

The deployed client asserts that the two bits agree on every round for which it
can observe both, and treats disagreement as a defect rather than as an input.

\paragraph{Reference implementation.} An open-source implementation of the
protocol, together with an independent reimplementation written against the
protocol's specification in a separate codebase with no shared code, is
available~\cite{fowrepo}. The two reproduce each other's flights, responses and
finishing bits byte for byte, so the specification and the deployed system
cannot silently diverge.

\section{Security analysis}
\label{sec:security}

The designed leakage is the bit, and, by padding, nothing about $|V|$.

\paragraph{The blinded set hides its contents --- stated with its hypothesis.}
Recovering anything about $V(a)$ from $S_A$ alone is the distinguishing problem
of \S\ref{sec:prelim}, requirement~2. The receiver, however, holds more than
$S_A$: it holds $Q_A = \alpha_A P_a$, and the response $\alpha_A Q_B$ hands it
one evaluation of $x \mapsto \alpha_A\,\HtoC(x)$ per round at a point of its
choice --- for a protocol-following receiver that point is its own committed
tile, which is exactly what \eqref{eq:finish} consumes. Set privacy against a
receiver is therefore a one-more/gap Diffie--Hellman-style claim
(\S\ref{sec:prelim}): the one evaluation tells the receiver whether its queried
point lies in the set, by design, and must help with nothing else. A receiver
that \emph{deviates} spends the same one query on an arbitrary point instead;
\S\ref{sec:oracle} measures what that is worth, and \S\ref{sec:qbind} is what
the deployed system does about it.

\paragraph{Unknown discrete logarithms are load-bearing.}
The blinding must be by a secret factor \emph{and} the elements being blinded
must have unknown discrete logarithms. If positions were instead encoded as
$P_t = h_t G$ with $h_t$ a public hash of $t$, then every element's discrete
logarithm would be public, and a receiver of one flight could strip the blinding:
each $h_v$ is public and invertible, so multiplying a received element by
$h_v^{-1}$ for candidate positions $v$ yields $\alpha G$ at the correct guess,
after which one further pass labels every element of the set. Over a game's small
universe of positions this is a few thousand scalar multiplications, so the
sender's set, and hence its position, is recovered from a single honest protocol
run. With a hash-to-group map the corresponding first step, obtaining $\alpha G$
from $\alpha P_t$ without knowing $\log_G P_t$, is the computational
Diffie--Hellman problem. This is also why brute force over a small universe buys
nothing: ``is this element position $t$?'' has no checkable relation without
breaking Diffie--Hellman.

\begin{remark}
This is not a hypothetical. A published reference implementation of ECDH-PSI
constructs its elements exactly this way~\cite{thomsonpsi}; we reported it
upstream, and it was fixed. It is the sharpest illustration we know of why the
hash-to-group map is not a stylistic detail.
\end{remark}

\paragraph{Cross-round unlinkability.}
$\alpha_r$ is fresh per round and derived from a secret that leaves the client
only after the epoch it governs is closed (\S\ref{sec:epochs}), so linking
$\alpha_r P_t$ to $\alpha_{r+1} P_t$, a stationary player, is again a
Diffie--Hellman distinguishing problem for an observer of the honest
transcript. The pad dummies are re-derived each round for exactly this reason:
elements recurring across rounds would become a linkable tag the moment any one
round's scalar went public, and a dispute filing publishes one scalar by design
(\S\ref{sec:derive}).

\paragraph{Reflection and replay.}
$B$ may echo one of $A$'s own set elements back as its $Q_B$, giving
$R_A = \alpha_A^2 P_t$. Exploiting the pair $(\alpha_A P_t,\, \alpha_A^2 P_t)$ is
the square Diffie--Hellman problem, for which no attack is known in the groups we
use. The bit that $B$ computes from a garbage $Q$ is garbage, which is
self-sabotage, and it becomes attributable evidence under \S\ref{sec:audit}.

\paragraph{Malformed elements.}
A malicious peer can substitute an arbitrary string for a group element. In our
x-only instantiation the ladder processes off-curve and low-order inputs without
special-casing; the worst case is a small-order element sent as $Q$, whose reply
$\alpha Q$ reveals $\alpha$ modulo the cofactor, a handful of bits out of the
scalar's full width, and useless against the distinguishing problem. The reply
never matches a cofactor-cleared set element, so the resulting bit is 0. Like
every malformed flight this is self-sabotage now and evidence later.

\subsection{The chosen-element oracle, and what it is worth}
\label{sec:oracle}

The deviation above deserves stating as an attack, because we measured it and
the measurement moved us. A responder cannot validate the query element --- a
responder able to tell what it is being asked would defeat the instrument ---
so a malicious player may aim its round's $Q$ at \emph{any} tile $t$: submit
$Q = P_t$, read the peer's response $\alpha_B P_t$, and test it for membership
of the peer's published set. Sight is symmetric, so the answer is ``are you
standing within sight of $t$'' --- a locator query about a place the asker is
nowhere near, answered truthfully by every peer, claiming nothing and filing
nothing.

Two prices attach to it even without the binding of \S\ref{sec:qbind}. The
asker forfeits its own warning for the round --- its instrument answers about
$t$, not about itself --- and none of the disclosure obligations move with the
question, so an asker genuinely in sight of its victim is convicted for staying
dark exactly like a deliberate suppressor. It is tempting to scope the leak by
those prices and by its rate limit of one bit per peer per round, and an
earlier internal review of this paper did. Simulation says otherwise: an
attacker that probes only when its belief set \emph{proves} no opponent can see
it takes zero conviction risk and still cuts a walking hider's median candidate
set from several hundred tiles (seventy-five for an objective-running policy)
to roughly ten, pinning it inside ten tiles on 49--71\% of rounds where an
honest hunter manages 16--26\%; only a sprinting policy, which also probes
least (a third of its rounds, against 39--78\% for the walking policies), keeps
its cover in the hundreds. One \emph{optimally chosen} bit per round keeps
pace with a one-step-per-round target indefinitely --- the fixed point of
dilate-then-halve at one step a round is a sixteen-tile candidate set, and the
simulated medians (8--11) sit at that order of magnitude; the same arithmetic
gives 64 tiles at two steps a round, which is the regime the mode-hiding below
forces. The lesson we record: counting a leak's bits
per round is not the same as simulating an adversary that spends them well.

The deployed response is two-fold. First, the wire no longer publishes the
declared movement mode below contact, so the attacker must dilate its belief
set by the maximum step count every round; the measured pinning rate collapses
from the 49--71\% above to under 7\% --- still a leak, no longer an
operationally decisive one. (The same zero-risk attacker, retooled for the
mode-hidden wire, probes on 28--87\% of its rounds --- the fraction the sampled
check of \S\ref{sec:qbind} prices.) Second, the query element is bound, retroactively
and at audit strength, to the committed position (\S\ref{sec:qbind}), which
turns the oracle from an unattributable capability into a self-incriminating
one. Neither is an in-round proof, and \S\ref{sec:limits} states the residue.

\paragraph{The accumulated-negative leak.}
Every ``no'' round tells a player that the opponent is outside its visible set,
and these constraints accumulate. This is the designed privacy and liveness
trade of the whole scheme rather than an oversight: any protocol that guarantees
requirement~2 of \S1 must leak at least the negatives, because a guaranteed
encounter is exactly the absence of a false negative. We measured the residual
anonymity set on our own maps against an exact Bayesian attacker granted every
advantage, and it remains large enough to play with, but the figure is a property
of a particular map, sight range and movement speed rather than of the protocol,
and it must be re-derived for any other setting.

\paragraph{What the exchange deliberately does not provide.}
First, no in-round proof that $V$ or $Q$ was computed from the committed
position: a malicious party can feed any set or $Q$ it likes, and the
surrounding system makes this pointless or attributable rather than impossible
(\S\ref{sec:audit}, \S\ref{sec:qbind}) --- with two bounded residues, an
over-reported set and the window before a chosen element's epoch is audited,
recorded honestly in \S\ref{sec:limits}.
Second, no fairness or liveness: a party can stall mid-exchange, and the
channel's existing timeout forfeits cover it. Third, it requires the predicate to
be symmetric (\S\ref{sec:predicate}).

\section{Attribution instead of zero-knowledge}
\label{sec:audit}

In-round, the referee performs no curve arithmetic over the exchange: its one
check is that a published set rebuilds the committed root, which is hashing,
and the correctness of the exchange's \emph{content} is never judged before an
audit or a dispute --- where curve work is bought deliberately and by measure,
one sampled check per audited epoch (\S\ref{sec:qbind}) and a fixed handful of
multiplications per dispute. The honest path never touches the chain --- a
channel consults it only to open, dispute and settle. Three facts combine into
the enforcement story.

\begin{enumerate}
\item \textbf{The honest transcript is a pure function.} By
\S\ref{sec:derive}, every byte party $i$ should send in round $r$ is
$T(\pseed_i, \petseed_i, r, \mathrm{position}_i(r), \mathrm{map})$, the seeds
those of the epoch containing $r$ (\S\ref{sec:epochs}): the scalar and the pad
from the PET seed, $V$ from the committed position and the public map, the
commitment salts from the path seed, and the order fixed by the sort.
\item \textbf{The audit reveals the preimages.} The pre-existing endgame audit
forces every player to reveal its epoch secrets and its path, verified against
those rounds' hash commitments, with closed epochs pinned by the mints that
ended them (\S\ref{sec:epochs}). Any judge can therefore recompute, after the
fact, the exact bytes an honest party would have sent in any round, and compare
them with what was sent. Each party publishes a commitment to the flight it sent
inside its own signed move, which makes a wrong byte non-repudiable. What the
judge may never do is \emph{regenerate} an honest flight from the seed --- a
hundred-odd curve derivations, some thirty times its per-call budget
--- so each check works from supplied or published bytes instead. The one
secret-derived quantity the adjudication needs, the disputer's round scalar,
rides the filing itself rather than being derived from any revealed seed; it
arrives unproven, and safely so, because every equality it enters compares it
against bytes its filer committed or published before the dispute existed, so a
forged scalar convicts the filer. On the
accused's side the evidence is one Merkle-authenticated element of the
committed set, and the recomputation a verdict costs the judge is one
hash-to-group map and two scalar multiplications per arm --- an amount chosen
against the fuel budget, since a first draft that checked twice as much did not
fit beside the audit in the same call. On the disputer's side a single
named element is not even possible: a dark opponent's position is revealed only
by the audit itself, after the evidence is frozen, and the test's own privacy
hides \emph{which} element matched --- so the disputer files its whole set, the
referee refuses the filing unless it rebuilds the committed root, and the judge
then locates the element itself by scanning bytes that are already bound.
Using the commitment the sender already publishes, rather than a signature
scheme inside the blob, is what keeps all of this affordable.
\item \textbf{Convictions never trust the bit.} ``In sight and stayed dark'' and
``disputed a sighting that was not there'' are both adjudicated by recomputing
$\sight$, Eq.~\eqref{eq:sight}, over positions bound by commitments made before
the round's information existed. The bit only tells honest clients when they are
obliged to act; the map convicts.
\end{enumerate}

So a proof system's job, \emph{prove your input was honest}, is replaced by
\emph{your input is a deterministic function of things you are already bound to
reveal, and lying about it is signed evidence}. The guarantee is delivered after
the fact instead of in-round, which a turn-based channel game can afford when
any player can force the audit, and the consensus-side cost is essentially zero.

\paragraph{The trust boundary, stated plainly.}
Attribution closes a gap that is open in both directions. A malicious opponent who
fabricates its flight can make an honest client's instrument read a phantom
``yes'', causing it to over-disclose or to file a dispute; and it can equally make
the instrument read a false ``no'', blinding the honest client so that it fails an
obligation it was never told it had. The second is the more dangerous, because it
is a winning strategy rather than a nuisance: blind an opponent every round, watch
it approach unaware, then escalate and dispute it, and a referee that judged the
contradiction on geometry alone would convict the blinded victim.

Both close only if the commitment to the flight PRECEDES the exchange, which is why
it rides the commitment move rather than the disclosure: the commit pass closes
before any set is published, so every participant's root is public before any
flight exists on the wire. The response needs no commitment of its own, since the
audit recomputes it from the revealed secrets and the peer's committed element;
publishing it is what makes a doctored response attributable. What the
adjudication then establishes, on the two positions the audit has verified, is
which side's committed bytes contradict the geometry: whether the disputer's own
set contained the element that would have informed its victim, and whether the
accused's set produced the hit the disputer reports. Convictions therefore follow
from committed data and never from either party's claimed bit, which is precisely
the quantity an attacker forges. A false accusation is not free, and neither is a
blinding carried out through the committed set.

One further property is required, and it is worth stating precisely because it is
a property of the transport rather than of the protocol: the bytes a party
DELIVERS must be the bytes it is later judged against. An earlier version of our
implementation exchanged the flights off-chain and bound the response only by
publication in a disclosure --- a disclosure ordered, by the dispute's own filing
rules, after the victim had already had to choose whether to reveal. A party could
therefore deliver a doctored response, or a doctored set, off-chain while
publishing the honest values, and the recipient could not distinguish doctored
from honest without the sender's secret scalar. What the audit can recompute is
what SHOULD have been sent, never what WAS; delivery is unprovable without
attribution on the delivered bytes, so no adjudication rule over off-chain
messages closes it. The deployed protocol therefore delivers nothing off-chain at
all. Both legs of the exchange are PUBLISHED as signed channel messages in a pass
of their own, after every commitment is fixed and before any disclosure: the
referee refuses a published set that does not rebuild the committed root, each
response is computed from the peer's committed element read out of the published
state --- so the exchange needs no delivered message anywhere --- and delivery and
publication coincide by construction, at which point the recomputation convicts a
doctored response exactly as it convicts a doctored set.

\subsection{Binding the blinded element}
\label{sec:qbind}

One input to the exchange is not covered by the story so far. The published set
is bound to a pre-committed root; positions are bound by hash commitments; but
the query element $Q$ was, in the terms above, merely \emph{published} ---
nothing tied its bytes to the tile committed beside it, and \S\ref{sec:oracle}
measures what an attacker does with that freedom. The binding cannot be checked
where $Q$ is used: the responder must not learn what it is being asked, and the
referee, which could know, cannot afford curve arithmetic in every round's
calls. So the binding is retroactive, in three parts.

\emph{Retention.} Each accepted commitment appends its $Q$ to an append-only
Merkle frontier in the player's state --- nine levels at our round cap, 288
bytes --- maintained beside the commitment chain and reset at each epoch
boundary.

\emph{A sampled check at the audit.} An audit reveal now also carries the
epoch's PET seed and one Merkle opening. The judge draws one round of the
audited epoch, from a hash over the running heads of every player's commitment
and disclosure chains, salted by an accumulator over the seeds already revealed
in the pass;
for any given prober those inputs are pinned by moves other players made after
its queries were already chosen, so it can neither steer nor predict the
draw. The judge then recomputes $\alpha_r$ from the revealed PET seed by
\eqref{eq:alpha}, recomputes $\alpha_r P_{t_r}$ from the revealed tile, and
verifies the opening against the frontier. A player that probed on $k$ of the
epoch's $n$ rounds is caught with probability $k/n$ --- 28--87\% of rounds for
the zero-risk strategies of \S\ref{sec:oracle} on the deployed wire, lowered
only by asking less. A mismatch \emph{rejects} the reveal rather than
convicting: nothing has pre-committed the PET seed, so a mismatch cannot be
attributed to a dishonest $Q$ rather than to a mis-supplied seed, and a
conviction may not hang on that ambiguity. Rejection forfeits the prober all
the same: its committed $Q$ opens under no seed it could supply, since the
probe point and the true tile point have no known scalar relation, so no
spelling of its reveal passes and the timeout ejects it.

\emph{Sealing, so closed epochs stay disputable.} The check above covers the
epoch the audit opens; without more, an epoch closed mid-match would launder
its probes. The canonical attack probes from the dark, walks into contact and
kills --- and the contact proof's own mint closes the probing epoch, after
which no later call holds its tiles. The mint therefore folds what it is about
to discard --- the frontier root, a hash of the epoch's verified tiles, the
spent PET seed, the boundary round and the round count --- into a per-player
accumulator that is never reset; the ratchet \eqref{eq:ratchet} has already
re-seeded behind it, so the disclosure compromises nothing after the boundary.
From the moment the epoch-closing reveal
publishes the spent PET seed, every opponent can verify every round of that
epoch off the wire for free; detection is certain, and only the conviction
needs a referee. A dedicated dispute move names one round of one closed epoch
and ships the accused's whole sealed history. The referee re-folds the
accumulator from zero --- a rebuild that misses the accused's standing value
convicts the \emph{filer}, since every input is public and the rebuild
deterministic --- then performs the one curve check at the named round,
convicting the accused on a mismatch and the filer on a match. Every
well-formed filing convicts somebody; a dispute that could cost nothing would
itself be a free probe. And conviction is sound here where the audit could only
reject, because every PET seed folded into the accumulator was sealed there by
the accused's own judge-accepted reveal: the accused authenticated the very
bytes that convict it.

\emph{Measured, on the deployed referee.} The audit reveal with its sampled
check is the system's worst call at 84.8\% of the per-call fuel budget; the
mint's fold adds about 1.4\% of the budget to the call that pays it; the
closed-epoch dispute, rebuilding a worst-case 383-round epoch and paying its
curve check in a move of its own, runs at 40.0\%; and the retained state fits a
four-player match in 7{,}865 of the platform's 8{,}192-byte state cap. The
residues --- one draw per audited epoch rather than a check per round, and the
kill-decided match that no audit pass ever reaches --- are stated in
\S\ref{sec:limits}.

\section{Related work}
\label{sec:related}

\paragraph{The Diffie--Hellman private-matching lineage.}
The double-blinding trick at the core of \S\ref{sec:protocol} is old and good.
Meadows~\cite{meadows86} used commutative blinding for cryptographic matchmaking
in 1986, and Huberman, Franklin and Hogg~\cite{hfh99} used the same structure for
private community matching in 1999. The modern literature knows it as ECDH-PSI:
Agrawal, Evfimievski and Srikant~\cite{aes03} formalised the
commutative-encryption approach for set intersection across private databases,
Jarecki and Liu~\cite{jl09,jl10} gave it security against malicious parties
under one-more-Diffie--Hellman-type assumptions, and De Cristofaro and
Tsudik~\cite{dt10} linear-complexity instantiations. (The other main PSI
family, oblivious polynomial evaluation over homomorphic encryption, begins
with Freedman, Nissim and Pinkas~\cite{fnp04}, who position it explicitly
against this lineage.) It runs
at planetary scale today, with Google's Password Checkup~\cite{checkup19} and
Apple's PSI system~\cite{applepsi} both descendants. Our protocol is this
classical construction \emph{reduced}: from set$\,\cap\,$set to membership of a
single element, evaluated in both directions, because the game needs one shared
bit and not the overlap.

\paragraph{Private set intersection for fog of war.}
The closest prior art to this work is OpenConflict~\cite{openconflict}. Motivated
by a practical memory-snooping map hack against real-time strategy games, it
distributes game state so that each client holds only what it is entitled to see,
using a private set intersection over grid cells, and pads the sets so that unit
counts do not leak. It reports 22~ms per synchronisation at the peak of a
professional StarCraft~II game, so the approach is fast enough for real time and
not merely for turn-based play. The padding idea in \S\ref{sec:derive} is theirs.
Two things differ here. OpenConflict assumes that players cannot lie about their
inputs, having neither a referee nor a commitment layer, whereas our setting is a
consensus engine in which a client may submit anything at all and the entire
question is what happens afterwards; and our predicate is a joint sighting test
adjudicated by a replicated referee, rather than a state-distribution mechanism
among clients who are trusted to be passive.

Applying that line of work to a blockchain game, with positions hash-committed
on chain and revealed at the end of the match for deterministic replay, was
sketched by Thomson~\cite{thomsonblog}, who also stated plainly that the case of
actively dishonest players remained unsolved. That is the gap
\S\ref{sec:audit} addresses: replaying positions and actions does not detect a
fabricated \emph{protocol} input, because nothing in such a scheme binds the
exchanged messages to the positions later revealed. Deriving every blinding
scalar from a committed seed is what closes it.

\paragraph{The implementation we started from.}
Our engineering entry point into this literature was Thomson's open-source
ECDH-PSI implementation~\cite{thomsonpsi}, a clear reference implementation of the
classical protocol with an encrypted-payload layer for revealing which elements
matched. Studying and prototyping against it shaped what we kept and what we
changed: we adopted an explicit hash-to-group map, for the reason in
\S\ref{sec:security}; dropped the payload-encryption layer entirely, since we
must never learn \emph{which} element matched and so have nothing to decrypt;
fixed the set cardinality by padding, since otherwise $|V|$ leaks for free; and,
most importantly, bound the inputs. Like most PSI reference implementations that
one leaves inputs unbound, so a malicious party can use the protocol as a
membership oracle over fabricated sets and query points, and our committed-seed
determinism, audit and query binding (\S\ref{sec:audit}, \S\ref{sec:qbind})
exist precisely to close that gap.

\paragraph{Private equality and proximity testing.}
Comparing secrets without revealing them goes back at least to Fagin, Naor and
Winkler~\cite{fnw96}. The closest prior art to our setting is Narayanan et al.'s
private proximity testing~\cite{narayanan11}, which reduces location proximity to
private equality of grid cells for social location-sharing. Our test differs in
the predicate, walkable-path distance plus line of sight rather than Euclidean
grid overlap; in being mandatory and bidirectional every round, since their users
choose when to query whereas our players are \emph{obliged} to act on the bit and
the obligation is enforced by consensus; and in the adversary handled, since
their protocols assume honest-but-curious phones whereas our clients are fully
malicious and handled by attribution.

\paragraph{Hidden information in adversarial games.}
Mental Poker~\cite{sra79} founded the field of playing hidden-information games
without a trusted dealer. The modern blockchain line is exemplified by Dark
Forest~\cite{darkforest}, which achieves fog of war on a public chain by
attaching a zk-SNARK to every move, proving it valid relative to hidden
coordinates. That is the ``prove your input was honest'' path we could not take,
since our referee's per-call budget is below the cost of a single Groth16
verification, and channel games settle on chains where verifier costs bind harder
still. The determinism-plus-audit construction of \S\ref{sec:audit} is our
substitute, and it is available to any turn-based game that already ends in a
full reveal. Note also the difference in what is proved: Dark Forest proves
unilateral facts about one player's hidden move, whereas a sighting is an
irreducibly \emph{joint} predicate over two players' secrets, which is what
forces a two-party protocol rather than a proof.

\paragraph{Infrastructure.}
Game channels are due to Kraft~\cite{kraft16}, deployed on Xaya, the successor of
Huntercoin, to our knowledge the first blockchain to run a multiplayer game world
fully on-chain, and whose scaling and privacy limits motivated this line of work.
Curve25519 is Bernstein's~\cite{bernstein06}, the Elligator~2 map is from
Bernstein, Hamburg, Krasnova and Lange~\cite{elligator}, the ladder and its test
vectors follow RFC~7748~\cite{rfc7748}, and RFC~9380~\cite{rfc9380} is the
general hash-to-curve framework.

\section{Limitations and future work}
\label{sec:limits}

\begin{enumerate}
\item \textbf{Attribution is deterrence, not prevention.} A participant who
fabricates a flight learns the answer it wanted before anything convicts it, so
the stake must exceed the value of the information. The window is bounded
differently per input: a fabricated \emph{set} is judged only inside a dispute's
adjudication, so its bound applies to fabrications that must surface in one,
while a fabricated \emph{query element} meets the sampled check with the
per-epoch probability of \S\ref{sec:qbind} and the closed-epoch dispute with
certain detection --- but each of those lands after the epoch, never in the
round the answer was read. Disputes are no longer terminal for the table at
more than two participants (\S\ref{sec:epochs}); at two they still are, since
the adjudication always convicts one of the pair. The demand-an-immediate-end
challenge and the round-cap audit remain terminal by design.
\item \textbf{An over-reported set is a bounded, conditionally priced probe.}
A participant may pad its published set with positions it cannot see, up to the
fixed pad size per round, forcing an exposure from whoever it guesses right
about. Padding is provable after the fact from the revealed seed, and a phantom
``yes'' hands its victim a \emph{winning} dispute --- the adjudication's
poisoning arm convicts the padder from its own committed bytes --- but the
victim must choose to spend that dispute, and no deployed rule prices the probe
when it does not. Pricing it unconditionally is future work.
\item \textbf{The query-element binding is retroactive and sampled.} Between
commission and audit a chosen element reads its answers freely
(\S\ref{sec:oracle}); the live check draws one round per audited epoch, so a
sparse prober trades catch probability against questions asked; and the dispute
needs a living filer with a disclosure pass to file in. At three or four
players a bystander files. At two players the canonical probe-then-hunt-then-kill
run ends the match at the kill, before any audit pass exists and with nobody
left to file, so it convicts nobody on-chain: the killer's own final reveal
publishes the epoch's PET seed and path, and its committed queries are already
on the wire, so the probing is publicly verifiable by anyone --- evidence
without a court. This dead-victim residual is the sharpest gap the binding
leaves. An in-round well-formedness proof for $Q$ --- a zero-knowledge
statement of $Q = \alpha_r P_{t_r}$ --- would close the oracle outright; it
does not fit the referee's budget today, and it is the sharpest candidate for
future work.
\item \textbf{A one-round rotation edge.} A suppressor whose disclosure happens
to be ordered after the honest player's can dodge a dispute for a single round of
fleeting edge-of-range contact.
\item \textbf{Privacy measured pairwise.} The deployed implementation runs the
exchange at up to four participants --- each round every participant publishes
one response per peer, and the dispute adjudication names its pair explicitly ---
at per-round cost linear in the number of pairs. The anonymity-set measurements
of \S\ref{sec:security}, however, are two-participant; more observers can only
shrink the hider's candidate set, and we have not yet quantified by how much.
\item \textbf{Cost.} The per-round curve arithmetic is a fraction of a
turn-based beat. Precomputing the position elements of \eqref{eq:tilepoint} per
map trims it by roughly a seventh at the median visible set (a quarter at its
largest); the pad of \eqref{eq:dummy} cannot be precomputed, by design
(\S\ref{sec:derive}), which is what removed the halving an earlier design
bought here. Much shorter beats would need batching or a faster field
implementation.
\item \textbf{The leak is a dial, not zero.} The accumulated-negative constraint
is inherent to guaranteed encounters, and we chose to measure the resulting
anonymity set rather than claim more than the construction gives.
\end{enumerate}

\section*{Acknowledgments}

Edward Thomson's open-source PSI implementation~\cite{thomsonpsi} was the
practical starting point of this work: a readable, working embodiment of the
classical protocol that let us prototype in days and decide what a consensus
setting had to change. The protocol itself stems from four decades of prior art
beginning with Meadows~\cite{meadows86}; our contribution is the adaptation, not
the trick.

\begin{thebibliography}{99}

\bibitem{meadows86}
C.~Meadows.
\newblock A more efficient cryptographic matchmaking protocol for use in the
absence of a continuously available third party.
\newblock In \emph{IEEE Symposium on Security and Privacy}, 1986.

\bibitem{hfh99}
B.~A. Huberman, M.~Franklin, and T.~Hogg.
\newblock Enhancing privacy and trust in electronic communities.
\newblock In \emph{ACM Conference on Electronic Commerce (EC)}, 1999.

\bibitem{fnw96}
R.~Fagin, M.~Naor, and P.~Winkler.
\newblock Comparing information without leaking it.
\newblock \emph{Communications of the ACM}, 39(5), 1996.

\bibitem{sra79}
A.~Shamir, R.~L. Rivest, and L.~M. Adleman.
\newblock Mental poker.
\newblock Technical Report MIT-LCS-TM-125, Massachusetts Institute of
Technology, 1979.

\bibitem{fnp04}
M.~J. Freedman, K.~Nissim, and B.~Pinkas.
\newblock Efficient private matching and set intersection.
\newblock In \emph{EUROCRYPT}, 2004.

\bibitem{aes03}
R.~Agrawal, A.~Evfimievski, and R.~Srikant.
\newblock Information sharing across private databases.
\newblock In \emph{ACM SIGMOD International Conference on Management of Data},
2003.

\bibitem{jl09}
S.~Jarecki and X.~Liu.
\newblock Efficient oblivious pseudorandom function with applications to
adaptive OT and secure computation of set intersection.
\newblock In \emph{Theory of Cryptography Conference (TCC)}, 2009.

\bibitem{jl10}
S.~Jarecki and X.~Liu.
\newblock Fast secure computation of set intersection.
\newblock In \emph{Security and Cryptography for Networks (SCN)}, 2010.

\bibitem{dt10}
E.~De~Cristofaro and G.~Tsudik.
\newblock Practical private set intersection protocols with linear
complexity.
\newblock In \emph{Financial Cryptography and Data Security}, 2010.

\bibitem{narayanan11}
A.~Narayanan, N.~Thiagarajan, M.~Lakhani, M.~Hamburg, and D.~Boneh.
\newblock Location privacy via private proximity testing.
\newblock In \emph{Network and Distributed System Security Symposium
(NDSS)}, 2011.

\bibitem{checkup19}
K.~Thomas et~al.
\newblock Protecting accounts from credential stuffing with password breach
alerting.
\newblock In \emph{USENIX Security Symposium}, 2019.

\bibitem{applepsi}
A.~Bhowmick, D.~Boneh, S.~Myers, K.~Talwar, and K.~Tarbe.
\newblock The Apple PSI system.
\newblock Apple Inc.\ technical report, 2021.

\bibitem{thomsonpsi}
E.~A. Thomson.
\newblock Private-Set-Intersection: an ECDH-based PSI implementation.
\newblock \url{https://github.com/EdwardAThomson/Private-Set-Intersection},
accessed July 2026.

\bibitem{openconflict}
E.~Bursztein, M.~Hamburg, J.~Lagarenne, and D.~Boneh.
\newblock OpenConflict: preventing real time map hacks in online games.
\newblock In \emph{IEEE Symposium on Security and Privacy}, 2011.

\bibitem{thomsonblog}
E.~A. Thomson.
\newblock Preventing cheaters in fog of war games.
\newblock \url{https://edward-thomson.medium.com/preventing-cheaters-in-fog-of-war-games-69f202fbe107},
June 2020.

\bibitem{fowrepo}
Xaya Developers.
\newblock fog-of-war: a blinded sighting test, reference implementation.
\newblock \url{https://github.com/xaya/fog-of-war}, 2026.

\bibitem{bernstein06}
D.~J. Bernstein.
\newblock Curve25519: new Diffie--Hellman speed records.
\newblock In \emph{Public Key Cryptography (PKC)}, 2006.

\bibitem{elligator}
D.~J. Bernstein, M.~Hamburg, A.~Krasnova, and T.~Lange.
\newblock Elligator: elliptic-curve points indistinguishable from uniform
random strings.
\newblock In \emph{ACM Conference on Computer and Communications Security
(CCS)}, 2013.

\bibitem{rfc7748}
A.~Langley, M.~Hamburg, and S.~Turner.
\newblock Elliptic curves for security.
\newblock RFC~7748, 2016.

\bibitem{rfc9380}
A.~Faz-Hern\'{a}ndez, S.~Scott, N.~Sullivan, R.~S. Wahby, and C.~A. Wood.
\newblock Hashing to elliptic curves.
\newblock RFC~9380, 2023.

\bibitem{kraft16}
D.~Kraft.
\newblock Game channels for trustless off-chain interactions in
decentralized virtual worlds.
\newblock \emph{Ledger}, 1:84--98, 2016.
\newblock \url{https://doi.org/10.5195/ledger.2016.15}.

\bibitem{darkforest}
Dark~Forest team.
\newblock Dark Forest: a zkSNARK-based incomplete-information game on
Ethereum.
\newblock \url{https://zkga.me}, 2020.

\end{thebibliography}

\end{document}
